\appendix

\section{Proof for Theorem 1}\label{sec:appendix:The1}

In this section, we provide the formal proof for the optimality of the \emph{Smallest-RLI-First} policy stated in Theorem 1.\subsection{Problem Setup and Definitions}We model the prefill process as a sequence of alternating communication and computation stages across layers $\ell = 1, \dots, N$.Let $L_{\text{curr}}(t)$ denote the layer index being computed or waiting to be computed at time $t$.For any communication flow $f$, let $L_{\text{target}}(f)$ be the layer where the data is consumed. The Relative Layer Index is defined as $\text{RLI}(f, t) = L_{\text{target}}(f) - L_{\text{curr}}(t)$.We adopt the assumptions stated in the theorem:\begin{enumerate}\item[(i)] \textbf{Ideal Computation:} Computation for layer $\ell$ starts immediately once its dependencies (flows with $L_{\text{target}} = \ell$) are met and runs for duration $C_\ell$ without interruption.\item[(ii)] \textbf{Fluid Model:} Communication flows are infinitely divisible, allowing preemption without overhead.\item[(iii)] \textbf{Dedicated Bandwidth:} We focus on a single link constraint (e.g., ingress or egress bottleneck).\end{enumerate}\subsection{Proof of Optimality}\textbf{Objective.} Minimizing the prefill makespan is equivalent to minimizing the completion time of the final computation layer. Since the total volume of computation is fixed, this is equivalent to minimizing the total duration of \emph{Execution Stalls} (intervals where the GPU is idle waiting for data).\textbf{Interval Classification.} We partition the timeline $[0, T]$ into two types of intervals based on the state of the system:\begin{itemize}\item \textbf{Stall Intervals ($\mathcal{I}_{\text{stall}}$):} Time periods where the GPU is idle because a flow $f$ with $\text{RLI}(f, t) = 0$ (i.e., targeting the current layer) is pending.\item \textbf{Overlap Intervals ($\mathcal{I}_{\text{overlap}}$):} Time periods where the GPU is actively performing computation. During these intervals, the link is available to transmit flows with $\text{RLI}(f, t) > 0$ (i.e., future layers).\end{itemize}\textbf{The Exchange Argument.}Suppose there exists an optimal schedule $\mathcal{S}^*$ that minimizes the makespan but violates the Smallest-RLI-First rule. This implies that at some time $t \in \mathcal{I}_{\text{stall}}$, the schedule assigns bandwidth to a flow $y$ with a larger RLI while a flow $x$ with a smaller RLI is pending.Specifically, since $t$ is a stall interval, there must be some pending flow $x$ with $\text{RLI}(x) = 0$ (blocking the current layer). The violation implies $\mathcal{S}^*$ serves flow $y$ where $\text{RLI}(y) > \text{RLI}(x) = 0$ during $[t, t+\delta]$.We construct a new schedule $\mathcal{S}'$ by applying an exchange:\begin{enumerate}\item \textbf{Swap:} In $\mathcal{S}'$, we assign the interval $[t, t+\delta]$ to flow $x$ instead of $y$.\item \textbf{Effect on Stall:} By serving $x$ earlier, the dependency for the current layer $L_{\text{curr}}$ is satisfied $\delta$ time units earlier (assuming $x$ was the bottleneck). This strictly reduces the duration of the current $\mathcal{I}_{\text{stall}}$ and advances the start of the next computation phase (and thus the next $\mathcal{I}_{\text{overlap}}$).\item \textbf{Feasibility of $y$:} Flow $y$ is displaced from $t$. However, since $\text{RLI}(y) > 0$, flow $y$ is not required until a future layer. The completion of $x$ triggers the start of computation $C_{L_{\text{curr}}}$, creating a new Overlap Interval. We can safely move the transmission of $y$ into this newly created Overlap Interval. Since $y$'s deadline is strictly later than $x$'s, this deferral does not violate any dependencies.\end{enumerate}\textbf{Conclusion.} The constructed schedule $\mathcal{S}'$ reduces the cumulative stall duration by transforming a portion of $\mathcal{I}_{\text{stall}}$ into $\mathcal{I}_{\text{overlap}}$. Repeating this exchange argument for all violations proves that the schedule strictly prioritizing the smallest RLI (serving $\text{RLI}=0$ flows immediately to minimize Stalls, and using Overlap intervals for $\text{RLI}>0$ flows) yields the minimum makespan. \qed

\section{Full algorithm for Robust inter-request scheduling} \label{sec:appendix:fullalg}

\begin{algorithm}[t!]
    \caption{Inter-request Scheduling}
    \small
    \label{alg:ETScheduling_Min}
    \begin{algorithmic}[1]
    \Procedure{InterScheduling}{$\mathcal{B}$, $M$, $B$}
        \FuncInput{$\mathcal{B}, M$: Batches and Traffic Matrices}
        \FuncInput{$B$: Global total drop budget}
        \FuncOutput{$\sigma, \mathcal{H}$}

        \State $S(\cdot) \gets 0$ \Comment{Interference from high-priority batches}
        \State $\mathcal{P} \gets \emptyset$; \quad $\mathcal{H} \gets \emptyset$

        \LineComment{Step 1: Sorting via Effective Deadline}
        \State $\sigma \gets \Call{SortByRED}{\mathcal{B}}$
        \State Precompute load vectors $L$ from $M$

        \For{$k \gets 1$ \textbf{to} $|\sigma|$}
            \State Let $i \gets \sigma[k]$
            \State $\mathcal{P} \gets \mathcal{P} \cup \mathcal{B}_i$ \Comment{Add to candidate pool}

            \State $\widehat{F}_i \gets \Call{EstFinishTime}{S, L_i}$
            \LineComment{Step 2: Feasibility Check}
            \While{$\widehat{F}_i > D^{Lo}_i$ \textbf{and} $|\mathcal{H}| < B$ \textbf{and} $\mathcal{P} \neq \emptyset$}
                \State $u^\star \gets \arg\max_u (S(u) + L_i(u))$ \Comment{Bottleneck}
                \LineComment{Step 3: selective pruning}
                \State $r^\star \gets \arg\max_{r \in \mathcal{P}} \text{Load of } r \text{ on } u^\star$

                \State $\mathcal{H} \gets \mathcal{H} \cup \{r^\star\}$; \quad $\mathcal{P} \gets \mathcal{P} \setminus \{r^\star\}$

                \If{$r^\star \in \mathcal{B}_i$} \Comment{Drop request in current batch}
                    \State $L_i(\cdot)\gets L_i(\cdot)-\ell_{r^\star}(\cdot)$
                \Else \Comment{Drop request in higher priority}
                    \State $S(\cdot)\gets S(\cdot)-\ell_{r^\star}(\cdot)$
                \EndIf

                \State $\widehat{F}_i \gets \Call{EstFinishTime}{S, L_i}$
            \EndWhile

            \State $S(\cdot) \gets S(\cdot) + L_i(\cdot)$ \Comment{Update}
        \EndFor

        \State \Return $\sigma, \mathcal{H}$
    \EndProcedure
    \end{algorithmic}
\end{algorithm}


The overall workflow of the scheduling algorithm is illustrated in Algorithm~\ref{alg:ETScheduling_Min}. In this section, we elaborate on the scheduling logic, providing the comprehensive mathematical formulation of priority assignment and further details on the practical implementation of latency estimation and enforcement.

This section details the robust scheduling algorithm designed to handle imprecise laxity and overload. The workflow operates in three phases triggered by batch arrival and departure events.

\parab{Step 1: Sorting via Robust Effective Deadline.}
To mitigate the "Piggyback Effect," the scheduler employs the Robust Effective Deadline ($RED$) metric defined in \S\ref{design:interreq}. The prioritization process executes in two stages. First, for each batch $\mathcal{B}$ with $n$ requests, the system sorts the requests internally by their deadlines $d_1 \le d_2 \le \dots \le d_n$. It then performs a linear scan to identify the partition point $k^* = \operatorname*{arg\,max}_{1 \le k < n} (d_{k+1} - d_k)$ that yields the maximal inter-request gap. This calculation dynamically separates the outliers (Tight Set) from the majority (Loose Set). Based on this partition, the scheduler computes the $RED$ score using the derived sub-batch deadlines and the outlier ratio $f$. Finally, all active batches are organized into a global priority queue sorted by ascending $RED$ values. This ensures that the dispatch order effectively filters out transient urgency spikes from isolated outliers while preserving the priority of genuinely urgent workloads.

\parab{Step 2: Worst-case Feasibility Check.} Following prioritization, we perform an admission control check to prevent the "Black Hole Effect." We estimate the completion time $\widehat{T}_i$ by strictly accumulating delays under a worst-case assumption (no overlap between batches). Computation latency is treated as deterministic, derived from the static computation graph of the Transformer model and offline profiling on the target hardware \cite{agrawal2024vidur,distserve,mooncake}. Communication latency is estimated by analyzing the batch's traffic matrix to identify the bottleneck port $p^*$ in the network fabric; the delay is calculated as the cumulative load on $p^*$ divided by its bandwidth. For dynamic architectures like Mixture-of-Experts (MoE) where token routing is runtime-dependent, we rely on historical routing statistics. Prior studies \cite{li2025optimizing} indicate that despite high noise in individual traffic matrices, the end-to-end latency prediction error typically remains within 20\%. This level of accuracy is sufficient for our scheduling granularity, allowing us to reliably enforce feasibility against the target deadline ($D^{Lo}_{\min}$).

\parab{Step 3: Selective Pruning and Soft Enforcement.} If a batch fails the feasibility check (i.e., $\widehat{T}_i > D^{Lo}_{\min}$), the system triggers a surgical pruning mechanism. We identify the specific requests within the batch that contribute the maximum load to the bottleneck port $p^*$ and iteratively remove them from the feasible set until the predicted latency meets the deadline. Crucially, this pruning employs a \emph{soft enforcement} strategy. Instead of immediately discarding the pruned requests, the scheduler demotes them to a "Scavenger" priority class within a dual-queue system. The Main Queue is serviced with strict priority, while the Scavenger Queue is processed opportunistically only when the network is detected to be idle or when actual runtime latencies are lower than the pessimistic estimates. This hybrid approach ensures system stability under worst-case predictions while maximizing throughput by reclaiming resources when congestion does not materialize.

\parab{Summary.} It is important to emphasize the design hierarchy within this workflow. The RED metric acts as the primary tie-breaker, governing the global dispatch order to maximize service-level objective (SLO) compliance under normal to moderate load. The feasibility check and selective pruning operate strictly as complementary safeguards, activated only during pathological overload to prevent resource exhaustion. This separation of concerns ensures that the scheduler remains stable and predictable—relying on robust prioritization for the vast majority of decisions—while retaining the capability to degrade gracefully only when extreme contention makes it unavoidable.
